\documentclass[11pt]{article}
\usepackage[margin=0.9in]{geometry}
\usepackage{amsmath,amssymb,amsthm,booktabs}
\usepackage[T1]{fontenc}
\usepackage{lmodern,microtype}
\usepackage[colorlinks=true,linkcolor=blue,urlcolor=blue]{hyperref}
\newtheorem{theorem}{Theorem}[section]
\newtheorem{lemma}[theorem]{Lemma}
\newtheorem{proposition}[theorem]{Proposition}
\newtheorem{corollary}[theorem]{Corollary}
\theoremstyle{remark}
\newtheorem{remark}[theorem]{Remark}
\newcommand{\norm}[1]{\left\|#1\right\|}
\title{Toward a four-dimensional Yang--Mills construction\\
\large Notebook 13: a uniform nonlinear constraint chart}
\author{Research working document}
\date{16 October 2026}

\begin{document}
% Notebook 14 body fragment; no standalone preamble.
% Requires amsmath/amssymb/amsthm and a one-argument \norm macro.
% Unique label prefix: nb-. Independently audited research proof.

\section{Ordinary reference backgrounds and their averaging operators}
\label{nb-section}

The corrected iteration at reference $1$ and the ordinary group
average have different flat linearizations. We relate the complete
configurations by an explicit gauge, then estimate the actual
moving-reference source operators at the ordinary backgrounds.
This argument does not equate those operators with the ordinary
derivative of the fixed-reference corrected iteration.

Fix odd $L\ge3$, $M=L^j$, $\eta=M^{-1}$ and
\begin{equation}\label{nb-scales}
 h_t=L^t\eta\quad(0\le t\le j),\qquad
 \sum_{s<t}h_s=\frac{h_t-\eta}{L-1}.
\end{equation}
We use compatible periodic lattices or the full lattice, globally
defined complex Lie-algebra bond fields, and physical volume in
every Hilbert space. All constants below are independent of $j$,
$M$ and compatible periodic volume. They may depend on fixed
$L$, dimension, color norm and the averaging prescription.
Matrix estimates use a submultiplicative norm; fixed norm
equivalences can be absorbed in the constants.

\subsection{The exact endpoint identity}
\label{nb-source}

Let $\mathsf M$ denote the ordinary group-valued average. For a
configuration $V$ near the identity put
\begin{equation}\label{nb-endpoint-map}
 \alpha(V;y)=L^{-d}\sum_{x\in B(y)}\frac1i
                      \log V(\Gamma_{y,x}),\qquad
 a(V;y)=e^{i\alpha(V;y)}.
\end{equation}
The corrected group map at reference $1$ is
\begin{equation}\label{nb-corrected-map}
 \mathsf C(V)(c)
   =a(V;c_-)^{-1}\mathsf M(V)(c)a(V;c_+).
\end{equation}
This is the specialization of deposited Raw I:24313--24325,
equivalently its explicit coordinate bridge I:47082--47086.
For a general reference the endpoint factor also carries an
adjoint by the average of that reference. It disappears here
because the reference is $1$.

Use the gauge convention
$V^g(c)=g(c_-)V(c)g(c_+)^{-1}$. Ordinary gauge covariance is
\begin{equation}\label{nb-gauge-covariance}
 \mathsf M(V^g)=\mathsf M(V)^{g|_{\rm coarse}}.
\end{equation}
The deposit records restriction to base points at
I:43018--43026, and expressly invokes complexified gauge
covariance from CMP98 (11) at I:82901--82904. In the common
identity neighborhood it also follows from similarity covariance
of the analytic products, logarithms and means. All gauges used
below stay inside that neighborhood.

These are source matches to the pinned deposited transcription.
The full original CMP98 paper has not been independently
reverified in this calculation. In particular, the bridge's
explicit incorporation of its omitted coordinate (120) is kept
as the source convention, rather than replaced by a guessed
coordinate formula.

\subsection{Uniform ordinary connection amplitudes}

Let $U=e^{i\eta A}$ and $a=\norm{A}_\infty$. Write
\begin{equation}\label{nb-two-trajectories}
 C_0=U,\quad C_{t+1}=\mathsf C(C_t),\qquad
 V_t=\mathsf M^t(U).
\end{equation}
The fixed-reference theorem of Notebook 13 gives a radius
$\rho>0$, physical corrected fields $A_t$ and a constant
$D_\infty$, such that for $a<\rho$
\begin{equation}\label{nb-corrected-input}
 C_t=e^{ih_tA_t},\qquad \norm{A_t}_\infty\le2a,
 \qquad A_{t+1}=\mathcal F_{h_t}(A_t),
\end{equation}
and
\begin{equation}\label{nb-one-step-derivative}
 \norm{D\mathcal F_h(B)}_{\infty\to\infty}
             \le1+D_\infty h\norm{B}_\infty.
\end{equation}
The derivative bound is used on a common smaller ball, which
contains the corrected trajectories and their straight-line
interpolations with translates.

\begin{theorem}\label{nb-trajectory-theorem}
There are $a_*>0$ and $C_V<\infty$, independent of refinement
and volume, such that if $\norm{A}_\infty<a_*$, every ordinary
reference $V_t$ has a holomorphic physical logarithm
\begin{equation}\label{nb-ordinary-amplitude}
 B_t=(ih_t)^{-1}\log V_t,\qquad
             \norm{B_t}_\infty\le C_V\norm{A}_\infty.
\end{equation}
Moreover there is an explicit vertex gauge $g_t$ with
$V_t=C_t^{g_t}$ and $g_t^{\pm1}-1=O(h_t\norm{A}_\infty)$.
All these maps read a bounded final-scale neighborhood.
\end{theorem}

\begin{proof}
Let $m_\Gamma$ be the largest number of oriented bonds in a
one-step contour. It is fixed when $L$ is fixed; the conservative
bound $m_\Gamma\le dL$ suffices. Decrease $a_*$ so that
$2m_\Gamma a_*\le\log(3/2)$. By
\eqref{nb-corrected-input},
\[
 \norm{C_t(\Gamma)-1}
      \le e^{2m_\Gamma h_ta}-1\le\tfrac12.
\]
The logarithm series and the bound
$e^s-1\le3s/2$ for $0\le s\le\log(3/2)$ give, with slack,
\begin{equation}\label{nb-alpha-bound}
 \norm{\alpha(C_t;y)}\le K_\Gamma h_ta,
                   \qquad K_\Gamma=8m_\Gamma.
\end{equation}
Set
\begin{equation}\label{nb-gauge-recursion}
 g_0=1,\qquad
 g_{t+1}(y)=g_t(y)a(C_t;y),\qquad y\in\Lambda_{t+1}.
\end{equation}
The lattices are nested, so the restriction of $g_t$ is defined.
If $V_t=C_t^{g_t}$, then \eqref{nb-gauge-covariance} and
\eqref{nb-corrected-map} imply
\[
 V_{t+1}(c)=g_t(c_-)a(C_t;c_-)C_{t+1}(c)
                       a(C_t;c_+)^{-1}g_t(c_+)^{-1}.
\]
Thus $V_t=C_t^{g_t}$ by induction, with exactly the product
order in \eqref{nb-gauge-recursion}.

The ordered exponential products and
\eqref{nb-scales}--\eqref{nb-alpha-bound} give
\begin{align}
 \norm{g_t^{\pm1}}
       &\le\exp\left(\frac{K_\Gamma a(h_t-\eta)}{L-1}\right),
       \label{nb-gauge-size}\\
 \norm{g_t^{\pm1}-1}
       &\le\exp\left(\frac{K_\Gamma a(h_t-\eta)}{L-1}\right)-1.
       \label{nb-gauge-difference}
\end{align}
The second bound follows by telescoping; no unitarity or
commutativity is required. Put
\begin{equation}\label{nb-explicit-constants}
 K_V=2+\frac{2K_\Gamma}{L-1},\qquad C_V=3K_V.
\end{equation}
The same product estimate for
$V_t(c)=g_t(c_-)e^{ih_tA_t(c)}g_t(c_+)^{-1}$ yields
\begin{equation}\label{nb-bond-near-identity}
 \norm{V_t(c)-1}\le e^{K_Vah_t}-1.
\end{equation}
Take also $K_Va_*\le\log(3/2)$. The analytic logarithm then
satisfies
\[
 \norm{\log V_t(c)}
    \le2(e^{K_Vah_t}-1)\le3K_Vah_t,
\]
which proves \eqref{nb-ordinary-amplitude}.

Further reduce $a_*$, if necessary, to place the right side of
\eqref{nb-bond-near-identity} in the fixed one-step analytic
neighborhood of $\mathsf M$. This also justifies the construction
inductively before each ordinary average is applied. Holomorphy
follows from the fixed-reference theorem and the explicit
products and logarithms. Each one-step dependency has physical
radius at most a fixed multiple of $h_t$, so the accumulated
radius is bounded by a multiple of the geometric sum
\eqref{nb-scales}. Alternatively, use the deposited two-final-block
support of the corrected bond iterate.
\end{proof}

\subsection{Physical first differences and curvature}

Supremum smallness alone does not give bounded physical
curvature. In a commuting direction take $A_2$ alternating
between $a$ and $-a$ at neighboring sites in direction $1$,
with all other components zero. The plaquette angle is of
order $\eta a$, hence the curvature divided by $\eta^2$ is
of order $a/\eta$. This obstruction is already present at $t=0$.

\begin{theorem}\label{nb-C1-theorem}
In addition to Theorem~\ref{nb-trajectory-theorem}, assume
\begin{equation}\label{nb-input-gradient}
 b=\max_\nu\norm{\nabla_{\eta,\nu}A}_\infty<\infty.
\end{equation}
Then, with fixed constants,
\begin{align}
 \max_\nu\norm{\nabla_{h_t,\nu}A_t}_\infty
       &\le e^{2D_\infty a/(L-1)}b,
                                      \label{nb-corrected-C1}\\
 \max_\nu\norm{\nabla_{h_t,\nu}B_t}_\infty
       &\le C e^{Ca}b,                  \label{nb-ordinary-C1}\\
 h_t^{-2}\norm{V_t(\partial p)-1}
       &\le C e^{Ca}(b+a^2).            \label{nb-curvature}
\end{align}
These estimates hold for real or complex fields on compatible
periodic lattices, with the same constants.
\end{theorem}

\begin{proof}
The fixed-block map $\mathcal F_h$ commutes with translations
by $Lh$ in each coordinate direction. For such a translation
$T$, integrate its derivative along the segment from $B$ to
$TB$. This segment remains in the same supremum ball. Equation
\eqref{nb-one-step-derivative} gives
\[
 \norm{\mathcal F_h(TB)-\mathcal F_h(B)}_\infty
  \le(1+D_\infty h\norm{B}_\infty)\norm{TB-B}_\infty.
\]
Telescoping $L$ fine steps shows
$\norm{TB-B}_\infty/(Lh)\le
\norm{\nabla_{h,\nu}B}_\infty$. Iterating, using
\eqref{nb-corrected-input} and \eqref{nb-scales}, proves
\eqref{nb-corrected-C1}. This requires only coarse translation
covariance, not invariance under an arbitrary fine translation.

On the same fixed coordinate ball, differentiation of the
finite contour formula gives
\begin{equation}\label{nb-contour-derivative}
 \norm{D_A\alpha(e^{ihA})}_{\infty\to\infty}
 +\norm{D_Aa(e^{ihA})}_{\infty\to\infty}\le C_\Gamma h.
\end{equation}
Indeed, each differentiated bond exponential contributes $h$;
the other factors have bounded product norm. The logarithm
derivative is bounded on $\norm{V-1}\le1/2$, as is the final
exponential derivative. No fine derivative is taken here.

For $s<t$, translation by $h_t$ preserves the scale-$s$ block
construction. By \eqref{nb-contour-derivative}, telescoping
$h_t/h_s$ input steps and using \eqref{nb-corrected-C1},
\begin{equation}\label{nb-endpoint-spatial-difference}
 \norm{a(C_s;y+h_te_\nu)-a(C_s;y)}
        \le C h_sh_t e^{2D_\infty a/(L-1)}b.
\end{equation}
The inverse factors obey the same estimate. Telescope the
ordered products defining $g_t$ at the two vertices, bound the
undifferentiated products by \eqref{nb-gauge-size}, and sum
$h_s$. This gives
\begin{equation}\label{nb-gauge-spatial-difference}
 \norm{g_t(y+h_te_\nu)^{\pm1}-g_t(y)^{\pm1}}
                         \le C e^{Ca}h_t^2b.
\end{equation}
The corresponding difference of $C_t(c)=e^{ih_tA_t(c)}$
is bounded by the same expression. The identity $V_t=C_t^{g_t}$
therefore bounds the difference of $V_t$ across $h_te_\nu$
by $Ce^{Ca}h_t^2b$. The logarithm is uniformly Lipschitz in
the half-unit ball, by integrating its derivative on the convex
segment. Dividing the logarithm difference by $h_t^2$ proves
\eqref{nb-ordinary-C1}. The case $t=0$ is the input estimate.

For curvature, first work with $C_t$. Pairing the two parallel
sides of a plaquette, the difference formula for exponentials
bounds the opposite-side differences by
$Ch_t^2\max_\nu\norm{\nabla_{h_t,\nu}A_t}_\infty$.
The remaining constant-side product is a group commutator,
bounded by $Ch_t^2a^2$. Thus
\[
 \norm{C_t(\partial p)-1}
 \le Ch_t^2\left(e^{2D_\infty a/(L-1)}b+a^2\right).
\]
The ordinary plaquette is conjugated by $g_t$ at its initial
vertex. Equation \eqref{nb-gauge-size} proves
\eqref{nb-curvature}. For real fields the gauges are unitary,
so a unitarily invariant norm avoids this conjugation loss.
Periodization preserves the translation identities and all
telescoping estimates, including on short compatible tori.
\end{proof}

\subsection{Actual bond averages in physical norms}

Let $Q_h(V)$ be the normalized linear coefficient of the
corrected relative increment at reference $V$, held fixed.
It is the one-step source operator of CMP98 (124), transcribed
at Raw I:24355--24369. The explicit bridge I:47082--47086
expresses it through finite local analytic operations. At $V=1$
it is the straight-string operator $Q_h^{\rm st}$.

For an operator from spacing $h$ to spacing $H$, define the
absolute physical Schur norm by
\begin{equation}\label{nb-av-norm}
 \norm{P}_{\rm av}
 =\max\left\{
       \sup_c\sum_b\norm{P_{cb}},\quad
       \sup_b(H/h)^d\sum_c\norm{P_{cb}}\right\}.
\end{equation}
This controls the bounded-field norm, the bounded-field norm
of the physical adjoint and the physical $L^2$ norm. It is
submultiplicative across scales: the volume ratios multiply.
For the flat normalized straight-string map both row sums
equal one. In fact each input bond occurs $L$ times among
the $L^d$ translated strings of length $L$, so its ordinary
column sum is $L^{-d}$ and its physical adjoint column sum
is one.

Finite-dimensional local analyticity and fixed-$L$ row and
column incidence give constants $r_Q,C_Q>0$ with
\begin{equation}\label{nb-local-Q-difference}
 \norm{Q_h(e^{iX})-Q_h^{\rm st}}_{\rm av}
             \le C_Q\norm{X}_\infty,
                 \qquad\norm{X}_\infty<r_Q.
\end{equation}
To verify this directly, differentiate each local coefficient
on a smaller closed complex polydisc and sum its finitely many
incidences. The physical column factor is $L^d$, independent
of $h$. The separate exact change-of-reference calculation
identifies the logarithm Jacobians and endpoint correction
that enter these coefficients; this shorter norm argument
uses their already explicit local analytic definition.

The actual source product uses the ordinary backgrounds:
\begin{equation}\label{nb-actual-product}
 Q_j(U)=Q_{h_{j-1}}(V_{j-1})\cdots Q_\eta(V_0).
\end{equation}
This is the composition recorded at Raw I:25147--25155 and
I:47046--47061. By Theorem~\ref{nb-trajectory-theorem},
$\norm{\log V_t}_\infty\le C_Vah_t$. Reduce $a_*$ so that
$C_Va_*<r_Q$. Submultiplicativity and an exact telescoping
of the difference of products prove
\begin{align}
 \norm{Q_j(U)}_{\rm av}
       &\le\exp\left(\frac{C_QC_Va}{L-1}\right),
                                      \label{nb-product-bound}\\
 \norm{Q_j(U)-Q_M^{\rm st}}_{\rm av}
       &\le\exp\left(\frac{C_QC_Va}{L-1}\right)-1.
                                      \label{nb-product-difference}
\end{align}
For a partial product ending at $h_t$, replace $(L-1)^{-1}$
in the exponent by $(h_t-\eta)/(L-1)$. These are actual source
operator estimates, obtained without renaming
$D\mathcal F^{(j)}(A)$ as $Q_j(U)$.

\subsection{Scalar averages and finite physical range}

The scalar average called $Q'_j$ in the deposit, and $Q_0$
in the homogeneous flat analysis, is
\begin{equation}\label{nb-scalar-average}
 (Q'_j(U)\lambda)(y)=L^{-jd}\sum_{x\in B^j(y)}
       \operatorname{Ad}_{U(\Gamma^{(j)}_{y,x})}\lambda(x).
\end{equation}
Raw I:23913--23920 states this formula. Its contour follows
successive block centers. The description at I:17798--17810
uses level-$s$ ordinary averaged configurations on the legs,
with at most $dL$ bonds per leg. The ordinary amplitude bound
and \eqref{nb-scales} therefore imply
\begin{equation}\label{nb-scalar-transport}
 \norm{U(\Gamma^{(j)}_{y,x})^{\pm1}}\le e^{Cah_j},
 \qquad
 \norm{\operatorname{Ad}_{U(\Gamma^{(j)}_{y,x})}-I}
                                      \le e^{Cah_j}-1.
\end{equation}
At a fixed level the blocks partition the fine vertices.
Each input scalar vertex occurs in one row with weight
$L^{-jd}$, canceled by the physical adjoint volume factor.
Consequently
\begin{equation}\label{nb-scalar-difference}
 \norm{Q'_j(U)}_{\rm av}\le e^{Cah_j},\qquad
 \norm{Q'_j(U)-Q'_j(1)}_{\rm av}\le e^{Cah_j}-1.
\end{equation}
This scalar operator is distinct from the bond average and
from the symbol $Q_0$ used for the main term of CMP98 (125).

Partition the spaces into compatible physical cells of side
one, with their physical Hilbert measures. The finite
multilevel reading radius is bounded independently of
refinement. For bond averages the deposit states the sharper
two-final-block support at I:23927--23930 and I:24332--24336.
If $R$ bounds the physical reading radius and $N_R$ the
number of possible neighboring cells, every cell block has
operator norm at most the global Hilbert norm. Hence
\begin{equation}\label{nb-cell-Schur}
 \norm{P}_{\kappa,\mathrm{cell}}
 :=\max\left\{
 \sup_x\sum_y e^{\kappa d(x,y)}\norm{P_{xy}}_{2\to2},
 \sup_y\sum_x e^{\kappa d(x,y)}\norm{P_{xy}}_{2\to2}\right\}
 \le N_Re^{\kappa R}\norm{P}_{2\to2}
\end{equation}
for such a finite-range $P$. Applying this to the differences
in \eqref{nb-product-difference} and
\eqref{nb-scalar-difference} gives uniform $O(a)$ bounds in
the physical-cell weighted Schur norm. Periodic distance and
fixed cover multiplicities give the same result on small
compatible tori.

\subsection{Local holomorphy and scope}

All constructions above are holomorphic on the same complex
ball $\norm{A}_\infty<a_*$. At finite regulators this follows
from products of finite-dimensional holomorphic maps. Their
uniform bounds and bounded physical reading ranges also give
the corresponding bounded-field local maps on the full
lattice. In particular $A\mapsto Q_j(e^{i\eta A})$ and
$A\mapsto Q'_j(e^{i\eta A})$ are holomorphic in the weighted
operator norm just defined, with a uniform bound on every
smaller ball. For example, if that bound is $C_0$ and an
argument lies at distance at least $r$ from the boundary of
the chosen ball, the multivariable Cauchy formula gives
\begin{equation}\label{nb-Cauchy}
 \norm{D^nQ_j(A)[K_1,\ldots,K_n]}_{\kappa,\mathrm{cell}}
 \le C_0 n!\left(\frac er\right)^n
                         \prod_{m=1}^n\norm{K_m}_\infty,
\end{equation}
and the same statement for $Q'_j$. One obtains it using
polydisc radii $r/(n\norm{K_m}_\infty)$ and
$n^n\le e^n n!$. Directions of zero norm are handled by
multilinearity. This is a bounded-background derivative
statement; no Hilbert norm is substituted for a background
direction without another estimate.

The result supplies uniform ordinary reference trajectories,
their physical $C^1$ and curvature bounds when the initial
first differences are bounded, and actual source averaging
operators close to their flat counterparts. It concerns a
globally trivialized homogeneous class. It does not establish
nonflat Hessian stability, the moving projected gauge inverse,
nonflat minimizing extensions, mixed retained-history domains,
strict background-support expansions or an interacting
Yang--Mills ultraviolet limit. The original-primary-paper
verification limitation stated in Section~\ref{nb-source}
remains in force.

\end{document}
